\section{Acondicionamiento, instrumentación y digitalización}
\label{sec:acondicionamiento-digitalizacion}

El frente analógico debe convertir una tensión diferencial pequeña en una secuencia que
aproveche el ADC, preserve la fase y se adquiera de forma determinista; la transmisión
se trata después.

Si $v_d(t)$ es la tensión diferencial de entrada,

\begin{equation}
\begin{aligned}
v_{\mathrm{ADC}}(t)={}&V_{\mathrm{CM}}+\left[h_{\mathrm{AFE}}(t)*v_d(t)\right]\\
&+v_{\mathrm{os}}(t)+n_{\mathrm{AFE}}(t),\\
x[n]={}&Q_{B,V_{\mathrm{span}}}\!\left\{v_{\mathrm{ADC}}(nT_s)\right\}.
\end{aligned}
\label{eq:afe-general}
\end{equation}

La ecuación reúne rechazo de modo común, ganancia programable, conformación espectral,
control del punto de operación, digitalización conocida y transferencia determinista a
memoria. El \emph{offset} acumulado puede agotar el rango aun sin señal.

\subsection{Compensador de polos desacoplados}

Tras evaluar varias arquitecturas se adoptó la forma $1-\mathrm{BP}$ propuesta e
implementada en \cite{Ma2023}. En ese trabajo, un pasa-banda Sallen--Key y un
amplificador diferencial restan la rama BP del camino directo.

La caída de \SI{27.7}{\decibel} a \SI{2}{\hertz} exige compensación antes del ADC si se
quiere observar contenido próximo a \SI{1}{\hertz}. Ma et al. muestran que conservar
$f_n$ y elevar el amortiguamiento efectivo extiende la respuesta aproximadamente plana a
aceleración entre

\begin{equation}
f_1\simeq\frac{f_n}{2\zeta_1},
\qquad
f_2\simeq2\zeta_1f_n.
\end{equation}

Por tanto, $\zeta_1$ se elige para llevar $f_1$ hasta la zona de interés. El valor de
$f_2$ no es un objetivo porque la aplicación no requiere extender la banda superior;
sólo debe conservarse la simetría logarítmica $f_1f_2=f_n^2$ alrededor de $f_n$.

Como los quiebres quedan muy separados, el $Q$ requerido es bajo: Sallen--Key o MFB no
aportan en este régimen y distribuyen cada polo entre más componentes, una desventaja
frente a tolerancias y valores disponibles.

Se desacoplaron los polos: uno en la entrada ($R_1$ serie $C_1$) y otro en la
realimentación ($R_2$ paralelo $C_2$) de un inversor. Cada uno depende de un solo par
$RC$. La inversión se reutiliza para combinar las ramas con un sumador inversor, sin
pareamiento, y su realimentación aporta el primer polo pasa-bajos.

\subsection{Ajuste y selección de componentes}

La cancelación depende de la proporción entre la rama directa y BP. Un potenciómetro en
serie con BP permite corregir esa proporción en cada unidad.

La red resultante se diseñó a partir del modelo dinámico del geófono,

\begin{equation}
H_{\mathrm{geo}}(s)=\frac{\omega_n s}
{s^2+2\zeta_0\omega_n s+\omega_n^2},
\end{equation}

donde $\zeta_0$ es el amortiguamiento nominal del SM-24. Se busca que la cadena formada
por el compensador y el geófono presente el amortiguamiento $\zeta_1$ requerido para
llevar el quiebre inferior a la zona de interés. Entonces

\begin{equation}
H_{\mathrm{comp}}(s)=
\frac{s^2+2\zeta_0\omega_n s+\omega_n^2}{s^2+2\zeta_1\omega_n s+\omega_n^2},
\end{equation}

y separar el término constante da la forma que se implementa en el circuito, un camino
directo menos un pasa-banda:

\begin{equation}
H_{\mathrm{comp}}(s)=1-\frac{2(\zeta_1-\zeta_0)\omega_n s}
{s^2+2\zeta_1\omega_n s+\omega_n^2}.
\label{eq:comp-afe}
\end{equation}

Se relevaron los pasivos disponibles y se incluyó el modelo no ideal del operacional del
PSoC: ganancia de continua de
\SI{90}{\decibel}, producto ganancia--ancho de banda de \SI{8}{\mega\hertz}, impedancia
de entrada de \SI{35}{\mega\ohm} y resistencia de salida de \SI{10}{\ohm}. Con ese modelo
se resolvió una búsqueda por evolución diferencial sobre la malla de valores comerciales
para minimizar el error respecto de la Ec.~\eqref{eq:comp-afe}. La búsqueda admitió
capacitores en paralelo.

Los quiebres necesarios para aproximar la banda buscada producen un amortiguamiento
efectivo muy alto: con los valores poblados se obtiene $f_0=\SI{10.21}{\hertz}$ y
$\zeta_1\simeq938$. En la zona central, el modelo compensado queda escalado por
$1/(2\zeta_1)\simeq5.3\times10^{-4}$, una atenuación de 1876 veces. La ganancia
se repartió para recuperar ese nivel sin concentrar todo el margen en un solo bloque:
la entrada instrumental
aporta $\times2$, el PGA es ajustable entre $\times1$ y $\times50$, el sumador aporta
$\times3.6$ y el pasa-bajos $\times5$. A máxima ganancia el producto es $\times1800$,
por lo que la señal del geófono llega al ADC con aproximadamente el \SI{96}{\percent}
de su escala previa a la compensación.

Como comprobación inicial se barrió la electrónica entre la salida monitoreada del PGA
y la entrada del ADC. La respuesta medida se encadenó en frecuencia con el modelo
nominal del SM-24 referido a aceleración para estimar la forma de la cadena completa.
En la realimentación del sumador se incorporó el polo real
$p_1=-2\pi(393)\,\mathrm{rad/s}$. La sección MFB se dimensionó con
$f_a=\SI{300}{\hertz}$ y $\zeta_a=0.691$; de
$p=-\zeta_a\omega_a\pm j\omega_a\sqrt{1-\zeta_a^2}$ resultan
$p_{2,3}\simeq2\pi(-207\pm j217)\,\mathrm{rad/s}$. El polo del sumador y este par
conjugado forman el antialias Bessel de tercer orden, elegido por su fase
aproximadamente lineal.

Las ganancias se mantuvieron conservadoras porque, al aumentarlas, cada paso del VDAC
también se amplifica: la calibración pierde resolución y la deriva, especialmente en el
sumador, se vuelve difícil de controlar. Como revisión pendiente se agregó una etapa
PGA de salida ajustable hasta $\times50$ y se sustituyeron las cuatro referencias VDAC
de \SIrange{0}{4.08}{\volt} por IDAC usados como DAC de tensión. Con la resistencia de
conversión propuesta, el ajuste se concentra aproximadamente entre
\SIrange{2.06}{3.02}{\volt}, en lugar de desperdiciar códigos fuera de la región útil.
La arquitectura permitiría otros $\times50$ de ganancia, pero aún requiere pruebas en
placa.

\begin{figure*}[p]
\centering
\includegraphics[width=0.97\textwidth,height=0.285\textheight,keepaspectratio]
{respuesta_geofono_afe_medida.png}
\caption{Cadena geófono--ADC: respuesta electrónica medida encadenada con el SM-24,
margen de ajuste y coherencia del barrido. La pérdida de coherencia próxima a
\SI{1}{\kilo\hertz} limita la interpretación superior; es una verificación de forma,
no una medición sísmica calibrada extremo a extremo.}
\label{fig:respuesta-geofono-afe}

\vspace{0.25em}
\textbf{(a) Ruta analógica empleada}\par
\includegraphics[width=0.99\textwidth,height=0.255\textheight,keepaspectratio]
{arquitectura_analogica_psoc.png}\par
\textbf{(b) Arquitectura digital}\par
\includegraphics[width=0.99\textwidth,height=0.255\textheight,keepaspectratio]
{arquitectura_digital_psoc.png}
\caption{Arquitectura utilizada en las mediciones: acondicionamiento, referencias,
AMux y ADC $\Delta\Sigma$; captura, filtrado, DMA, memoria y temporización.}
\label{fig:topdesign-analogico}
\end{figure*}

\begin{figure*}[!t]
\centering
\includegraphics[width=0.99\textwidth,height=0.29\textheight,keepaspectratio]
{propuesta_idac_pgaout_nueva.png}
\caption{Nueva propuesta para el frente GEO: PGA de salida y cuatro IDAC concentrados
en la región útil de ajuste. Aporta hasta $\times50$ adicionales y aún no fue empleada
en las mediciones.}
\label{fig:propuesta-idac-pgaout}
\end{figure*}

\subsection{Acople DC y calibración}

Al extender la compensación hacia \SI{1}{\hertz}, el \emph{offset} y la deriva térmica,
en especial del sumador, consumieron el margen del ADC aun sin señal.

Un multiplexor observa cada nodo con el mismo ADC y cada etapa posee referencia
programable. La calibración \emph{foreground} recorre PGA, BP, sumador y pasa-bajos; un
FIR de 128 taps estima el residuo y un PI con banda muerta, saturación y
\emph{anti-windup} corrige la referencia. El orden evita propagar el residuo aguas abajo.
Los códigos se guardan en EEPROM por ganancia, con CRC. Durante la captura el lazo se
desactiva y los códigos quedan fijos.

\subsection{Digitalización y memoria}

El ADC $\Delta\Sigma$ del PSoC~5LP entrega 18 bits nominales a
$F_{s,0}=\SI{2604}{SPS}$ y cuatro rangos bipolares entre $\pm\SI{0.512}{\volt}$ y
$\pm\SI{2.5}{\volt}$. PGA y rango se eligen juntos para usar el menor rango sin saturar.
El DFB filtra por FIR; dos DMA trasladan crudo y filtrado a RAM; y
un bloque de control implementado en Verilog/UDB coordina inicio, ADC, conteo y fin.
El software queda para calibración, persistencia y comandos.

\subsection{Estado de verificación}

Queda medir magnitud y fase extremo a extremo para cada ganancia y rango, ruido referido
a entrada y bits efectivos, dispersión entre placas, deriva térmica y repetibilidad de
la calibración tras migrar las referencias VDAC a IDAC.

\begin{figure}[!tb]
\centering
\includegraphics[width=\columnwidth,keepaspectratio]
{flujo_calibracion_horizontal_usuario.png}
\caption{Máquina de estados de la calibración \emph{foreground}: valida los códigos,
ajusta PGA, BP, sumador y pasa-bajos, y restaura la adquisición con códigos fijos.}
\label{fig:maquina-estados-calibracion}
\end{figure}
